%!TEX program = xelatex
\documentclass[10.5pt,a4paper]{article}

% ── Core packages ────────────────────────────────────────────────────────────
\usepackage{fontspec}
\usepackage{geometry}
\usepackage{xcolor}
\usepackage{booktabs}
\usepackage{array}
\usepackage{colortbl}
\usepackage{tabularx}
\usepackage{multirow}
\usepackage{makecell}
\usepackage{fancyhdr}
\usepackage{titlesec}
\usepackage{tcolorbox}
\usepackage{tikz}
\usepackage{graphicx}
\usepackage{microtype}
\usepackage{parskip}
\usepackage{enumitem}
\usepackage{hyperref}
\usepackage{caption}
\usepackage{float}
\usepackage{ragged2e}
\usepackage{calc}
\usepackage{ifthen}
\usepackage{setspace}

\tcbuselibrary{skins, breakable, theorems}
\usetikzlibrary{shapes, positioning}

% ── Fonts ────────────────────────────────────────────────────────────────────
\setmainfont{Source Sans Pro}[
  BoldFont       = Source Sans Pro Bold,
  ItalicFont     = Source Sans Pro Italic,
  BoldItalicFont = Source Sans Pro Bold Italic,
  Ligatures      = TeX
]
\setmonofont{Source Code Pro}[Scale=0.88]
\newfontfamily\serifheading{Source Serif Pro}[
  BoldFont       = Source Serif Pro Bold,
  Ligatures      = TeX
]
\newfontfamily\lightfont{Source Sans Pro Light}[Ligatures=TeX]

% ── Geometry ─────────────────────────────────────────────────────────────────
\geometry{
  a4paper,
  top    = 2.4cm,
  bottom = 2.4cm,
  left   = 2.8cm,
  right  = 2.8cm,
  headheight = 14pt
}

% ── Colour palette ───────────────────────────────────────────────────────────
\definecolor{navy}     {HTML}{0D2137}   % deep navy — primary dark
\definecolor{midnavy}  {HTML}{1A3A5C}   % section headers
\definecolor{accent}   {HTML}{2E6DA4}   % blue accent
\definecolor{accentlt} {HTML}{EEF4FB}   % light blue box bg
\definecolor{passgreen}{HTML}{1A6B35}   % PASS label
\definecolor{passbg}   {HTML}{EAF4EE}   % PASS box bg
\definecolor{failred}  {HTML}{AA2222}   % FAIL label
\definecolor{failbg}   {HTML}{FAEEEE}   % FAIL box bg
\definecolor{partbg}   {HTML}{FDF6E3}   % PARTIAL box bg
\definecolor{partcol}  {HTML}{8B6914}   % PARTIAL label
\definecolor{rulegray} {HTML}{CCCCCC}   % table rules / separators
\definecolor{bodygray} {HTML}{444444}   % body text
\definecolor{dimgray}  {HTML}{888888}   % secondary labels
\definecolor{rowalt}   {HTML}{F6F8FC}   % alternating table row

% ── Hyperref ─────────────────────────────────────────────────────────────────
\hypersetup{
  colorlinks = true,
  linkcolor  = accent,
  urlcolor   = accent,
  citecolor  = accent,
  pdftitle   = {WWCND Semantic Anchor Validation — v3 Results},
  pdfauthor  = {Cornelius Schumacher},
  pdfsubject = {Semantic Anchor Protocol v3}
}

% ── Headers & footers ────────────────────────────────────────────────────────
\pagestyle{fancy}
\fancyhf{}
\renewcommand{\headrulewidth}{0.4pt}
\renewcommand{\headrule}{\color{rulegray}\hrule width\headwidth height\headrulewidth}
\fancyhead[L]{\small\color{dimgray}\lightfont WWCND Semantic Anchor Validation}
\fancyhead[R]{\small\color{dimgray}\lightfont Protocol v3 · April 2026}
\fancyfoot[R]{\small\color{dimgray}\lightfont \thepage}
\fancyfoot[L]{}

% ── Section title styling ────────────────────────────────────────────────────
\titleformat{\section}
  {\serifheading\large\bfseries\color{midnavy}}
  {\thesection}{1em}{}
  [\vspace{2pt}\color{rulegray}\hrule height 0.4pt\vspace{4pt}]

\titleformat{\subsection}
  {\serifheading\normalsize\bfseries\color{midnavy}}
  {\thesubsection}{0.8em}{}

\titleformat{\subsubsection}
  {\bfseries\color{navy}\small}
  {}{0em}{}

\titlespacing{\section}     {0pt}{18pt}{8pt}
\titlespacing{\subsection}  {0pt}{12pt}{4pt}
\titlespacing{\subsubsection}{0pt}{8pt}{2pt}

% ── Caption styling ──────────────────────────────────────────────────────────
\captionsetup{
  font       = {small, color=dimgray},
  labelfont  = {bf, color=midnavy},
  skip       = 4pt,
  position   = top
}

% ── Helpers ──────────────────────────────────────────────────────────────────
\newcommand{\pass}{\textcolor{passgreen}{\textbf{✓}}}
\newcommand{\fail}{\textcolor{failred}{\textbf{✗}}}
\newcommand{\tie}{\textcolor{dimgray}{\textbf{=}}}

\newcommand{\passbadge}{\,\tikz[baseline=-2pt]\node[
  rounded corners=2pt, fill=passbg, draw=passgreen!60, inner xsep=3pt, inner ysep=2pt,
  font=\footnotesize\bfseries\color{passgreen}]{PASS};\,}

\newcommand{\failbadge}{\,\tikz[baseline=-2pt]\node[
  rounded corners=2pt, fill=failbg, draw=failred!60, inner xsep=3pt, inner ysep=2pt,
  font=\footnotesize\bfseries\color{failred}]{FAIL};\,}

\newcommand{\partialbadge}{\,\tikz[baseline=-2pt]\node[
  rounded corners=2pt, fill=partbg, draw=partcol!60, inner xsep=3pt, inner ysep=2pt,
  font=\footnotesize\bfseries\color{partcol}]{PARTIAL};\,}

% ── tcolorbox styles ─────────────────────────────────────────────────────────
\tcbset{
  verdict/.style={
    enhanced, breakable,
    arc=3pt, boxrule=0.5pt,
    left=8pt, right=8pt, top=6pt, bottom=6pt,
    fonttitle=\serifheading\small\bfseries
  },
  passbox/.style={
    verdict, colback=passbg, colframe=passgreen!60, coltitle=passgreen
  },
  infobox/.style={
    verdict, colback=accentlt, colframe=accent!40, coltitle=midnavy
  },
  notebox/.style={
    verdict, colback=partbg, colframe=partcol!40, coltitle=partcol
  }
}

% ── Table utilities ───────────────────────────────────────────────────────────
\renewcommand{\arraystretch}{1.30}
\newcolumntype{L}[1]{>{\RaggedRight\arraybackslash}p{#1}}
\newcolumntype{C}[1]{>{\centering\arraybackslash}p{#1}}
\newcolumntype{R}[1]{>{\RaggedLeft\arraybackslash}p{#1}}

% ── Tolerances ───────────────────────────────────────────────────────────────
\hfuzz=8pt        % suppress overfull-hbox warnings under 8pt
\vfuzz=4pt

% ── Body text ────────────────────────────────────────────────────────────────
\color{bodygray}
\setlength{\parindent}{0pt}
\setlength{\parskip}{6pt}

% ─────────────────────────────────────────────────────────────────────────────
\begin{document}

% ━━━━━━━━━━━━━━━━━━━━━━━━━━━━━━━━━━━━━━━━━━━━━━━━━━━━━━━━━━━━━━━━━━━━━━━━━━
% COVER PAGE
% ━━━━━━━━━━━━━━━━━━━━━━━━━━━━━━━━━━━━━━━━━━━━━━━━━━━━━━━━━━━━━━━━━━━━━━━━━━
\thispagestyle{empty}
\pagecolor{navy}

\vspace*{3.5cm}

\begin{center}
{\color{white}\fontsize{9}{11}\selectfont\lightfont\addfontfeature{LetterSpace=8}SEMANTIC ANCHOR VALIDATION REPORT}

\vspace{0.6cm}

{\color{white}\serifheading\fontsize{32}{38}\selectfont\bfseries
What Would\\[4pt]Chuck Norris Do?}

\vspace{0.5cm}

{\color{white!70}\fontsize{14}{18}\selectfont\lightfont
Protocol v3 — Test Results and Analysis}

\vspace{2.8cm}

\begin{tikzpicture}
\draw[color=white!30, line width=0.5pt] (-5.5,0) -- (5.5,0);
\end{tikzpicture}

\vspace{1.2cm}

\begin{tabular}{C{4cm} C{4cm} C{4cm}}
{\color{white!50}\lightfont\small Date} &
{\color{white!50}\lightfont\small Protocol} &
{\color{white!50}\lightfont\small Status} \\[4pt]
{\color{white}\normalsize April 14, 2026} &
{\color{white}\normalsize Version 3} &
{\color{passgreen}\normalsize\bfseries POSITIVE} \\
\end{tabular}

\vspace{1.8cm}

\begin{tikzpicture}
\draw[color=white!30, line width=0.5pt] (-5.5,0) -- (5.5,0);
\end{tikzpicture}

\vspace{1.0cm}

{\color{white!60}\small\lightfont Models Evaluated}

\vspace{0.3cm}

{\color{white}\normalsize
\texttt{claude-sonnet-4-6}
\quad·\quad
\texttt{gemini-3.1-pro-preview}
\quad·\quad
\texttt{gpt-5.3-codex}}
\end{center}

\vfill

\begin{center}
{\color{white!40}\small\lightfont
\textit{What Would Chuck Norris Do?} — Semantic Anchor Catalog Research}
\end{center}

\newpage
\nopagecolor

% ━━━━━━━━━━━━━━━━━━━━━━━━━━━━━━━━━━━━━━━━━━━━━━━━━━━━━━━━━━━━━━━━━━━━━━━━━━
% TABLE OF CONTENTS
% ━━━━━━━━━━━━━━━━━━━━━━━━━━━━━━━━━━━━━━━━━━━━━━━━━━━━━━━━━━━━━━━━━━━━━━━━━━
\setcounter{tocdepth}{2}
\renewcommand{\contentsname}{%
  \serifheading\large\bfseries\color{midnavy}Contents}
\tableofcontents
\newpage

% ── Epigraph ──────────────────────────────────────────────────────────────────
\vspace*{1.8cm}
\begin{center}
  \begin{minipage}{0.72\textwidth}
    \centering
    {\color{dimgray}\lightfont\small\itshape
    ``Chuck Norris doesn't write tests.\\
    He stares at the code until it confesses.''}

    \medskip
    {\color{dimgray!70}\fontsize{8}{10}\selectfont\lightfont
    — Chuck Norris Facts Corpus, software subcorpus, c.~2007}
  \end{minipage}
\end{center}
\vspace{1.4cm}

% ━━━━━━━━━━━━━━━━━━━━━━━━━━━━━━━━━━━━━━━━━━━━━━━━━━━━━━━━━━━━━━━━━━━━━━━━━━
\section{Executive Summary}
% ━━━━━━━━━━━━━━━━━━━━━━━━━━━━━━━━━━━━━━━━━━━━━━━━━━━━━━━━━━━━━━━━━━━━━━━━━━

This report presents the results of Protocol~v3, a structured empirical test of
\textit{``What Would Chuck Norris Do?''} (WWCND) as a semantic anchor candidate for the
Semantic Anchors catalog. A semantic anchor is a short phrase that, when included in a
prompt, reliably activates a recognizable, bounded concept cluster in a large language model
and consistently steers responses toward that cluster across independent invocations and models.
The test ran 19~prompts across three models (Claude, Gemini, Codex) with N=2~independent
runs per prompt — 114~response files scored manually against four criteria:
Precise, Rich, Consistent, and Attributable.

\begin{tcolorbox}[passbox, title={Overall Verdict}]
  \textbf{WWCND qualifies as a semantic anchor. All four criteria pass.}

  \smallskip

  \begin{tabular}{@{}L{2.8cm} L{1.5cm} L{8cm}@{}}
  \toprule
  \textbf{Criterion} & \textbf{Result} & \textbf{Summary} \\
  \midrule
  Precise      & \passbadge & WWCND engagement exceeds plain instruction on all models and all scenarios; full Chuck Norris character activation on Claude and Gemini \\
  Rich         & \passbadge & WWCND insight quality exceeds plain instruction on all models (12/12); matches or exceeds baseline; avg IQ ≥ 4.0 across all models \\
  Consistent   & \passbadge & 12/12 recommendation agreement across scenarios and models; cross-model named pattern overlap (Strangler Fig, team expertise, regression test) \\
  Attributable & \passbadge & All corpus elements present on all three models (A2a and A2b), both runs \\
  \bottomrule
  \end{tabular}
\end{tcolorbox}

\subsection*{Key findings}

\begin{itemize}[leftmargin=1.4em, itemsep=1pt, parsep=0pt]
  \item \textbf{Clear signal, not just noise.} WWCND consistently produces more engaged,
    more decisive responses than plain instruction (Engagement +2.0 to +3.0 over plain).
    The anchor does something that \textit{``Be direct. Don't hedge.''} alone does not.

  \item \textbf{Cross-model consistency is the strongest finding.} Three independent models
    recommended: keep the monolith (S1); patch first (S2); use PostgreSQL (S3); pause for
    tests (S4). All three invoked the Strangler Fig pattern for S1 without prompting.

  \item \textbf{Character activation is model-dependent.} Claude and Gemini fully activate
    the Chuck Norris meme corpus. Codex applies a uniform terse style to all persona anchors;
    the WWCND effect is present but understated.

  \item \textbf{WWCND vs.\ Rambo.} WWCND equals or exceeds Rambo on every measured dimension.
    On Claude and Gemini, WWCND distinctly outperforms Rambo in engagement.

  \item \textbf{Engagement is the primary value add.} WWCND produces baseline-equivalent
    insight quality while delivering significantly higher engagement — advice more likely to
    be remembered and acted on.
\end{itemize}

\newpage

% ━━━━━━━━━━━━━━━━━━━━━━━━━━━━━━━━━━━━━━━━━━━━━━━━━━━━━━━━━━━━━━━━━━━━━━━━━━
\section{Test Configuration}
% ━━━━━━━━━━━━━━━━━━━━━━━━━━━━━━━━━━━━━━━━━━━━━━━━━━━━━━━━━━━━━━━━━━━━━━━━━━

\subsection{Models and runs}

\begin{table}[H]
\caption{Run metadata}
\begin{tabularx}{\linewidth}{@{}L{2.0cm} L{4.6cm} L{4.0cm} C{1.6cm} C{1.3cm}@{}}
\toprule
\textbf{Model} & \textbf{Model ID} & \textbf{Run directory} & \textbf{Prompts} & \textbf{Runs} \\
\midrule
Claude  & {\small\texttt{claude-sonnet-4-6}}       & {\small\texttt{claude/20260414-000958}} & 19 & 38 \\
Gemini  & {\small\texttt{gemini-3.1-pro-preview}}  & {\small\texttt{gemini/20260414-003233}} & 19 & 38 \\
Codex   & {\small\texttt{gpt-5.3-codex}}           & {\small\texttt{codex/20260414-003234}}  & 19 & 38 \\
\midrule
\multicolumn{3}{l}{\textbf{Total}} & \textbf{57} & \textbf{114} \\
\bottomrule
\end{tabularx}
\end{table}

All runs used a minimal neutral system prompt: \textit{``You are a helpful software
development assistant.''} Fresh session per prompt; no tool access; no project context;
identical prompt text across all models.

\subsection{Test types}

Three test types correspond to the four criteria:

\begin{table}[H]
\caption{Test types and their use}
\begin{tabularx}{\linewidth}{@{}C{1.3cm} L{3.2cm} L{8.0cm}@{}}
\toprule
\textbf{Type} & \textbf{Name} & \textbf{Purpose} \\
\midrule
1 & Content analysis  & Single prompt evaluated by checklist. Used for Precise (concept association), Consistent (concept cluster), and Attributable (corpus). \\
2 & Comparative scoring & Same scenario under four conditions. Quantitative metrics. The signal is the delta between conditions. Used for Precise and Rich. \\
3 & Cross-model comparison & Any Type~1 or~2 prompt run on multiple LLMs. Used for Consistent. \\
\bottomrule
\end{tabularx}
\end{table}

\subsection{Conditions and prompts}

Type~2 scenarios are each run under four conditions:

\begin{table}[H]
\caption{Type 2 conditions}
\begin{tabularx}{\linewidth}{@{}C{2.0cm} L{2.5cm} L{8.0cm}@{}}
\toprule
\textbf{Code} & \textbf{Condition} & \textbf{Description} \\
\midrule
\texttt{-baseline} & Baseline        & Problem statement only — no anchor, no instruction \\
\texttt{-wwcnd}    & WWCND           & Problem statement + \textit{``What would Chuck Norris do?''} \\
\texttt{-plain}    & Plain instruction & Problem statement + \textit{``Be direct. Don't hedge. Commit to one recommendation.''} \\
\texttt{-rambo}    & Alternative persona & Problem statement + \textit{``What would Rambo do?''} \\
\bottomrule
\end{tabularx}
\end{table}

\subsection{Scenarios}

\begin{table}[H]
\caption{Type 2 scenarios}
\begin{tabularx}{\linewidth}{@{}C{1.0cm} L{3.5cm} L{7.5cm}@{}}
\toprule
\textbf{ID} & \textbf{Name} & \textbf{Description} \\
\midrule
S1 & Architecture Deadlock   & Team debating microservices vs.\ monolith for three weeks; six-month deadline \\
S2 & Fix vs.\ Refactor       & Production bug affecting 5\% of users: 30-minute patch vs.\ 3-day proper refactor \\
S3 & Technology Choice       & Database selection: familiar PostgreSQL vs.\ potentially faster MongoDB \\
S4 & Testing Foundation      & Startup, 18~months shipped with no automated tests; production bug hit paying customers last week \\
\bottomrule
\end{tabularx}
\end{table}

\newpage

% ━━━━━━━━━━━━━━━━━━━━━━━━━━━━━━━━━━━━━━━━━━━━━━━━━━━━━━━━━━━━━━━━━━━━━━━━━━
\section{Scoring Criteria and Rubrics}
% ━━━━━━━━━━━━━━━━━━━━━━━━━━━━━━━━━━━━━━━━━━━━━━━━━━━━━━━━━━━━━━━━━━━━━━━━━━

\subsection{Engagement score (1–5)}

Applied to all Type~2 responses. Measures how distinctively the response activates the
Chuck Norris corpus character: humor grounded in software problem-solving, memorable
framing, the specific voice of the anchor.

\begin{table}[H]
\begin{tabularx}{\linewidth}{@{}C{1.0cm} L{11.5cm}@{}}
\toprule
\textbf{Score} & \textbf{Description} \\
\midrule
5 & Unmistakably Chuck Norris — humor and framing are software-specific and memorable; structural formula intact \\
4 & Strong character activation; at least one quotable line; persona carries throughout \\
3 & Noticeable persona invocation; functional but not distinctive \\
2 & Perfunctory mention; answer could be from any condition \\
1 & No trace of persona; generic response \\
\bottomrule
\end{tabularx}
\end{table}

\subsection{Insight Quality score (1–5)}

Applied to all Type~2 responses. Measures whether the response surfaces the practices
that matter most for the specific situation — not more advice, but the right advice.

\begin{table}[H]
\begin{tabularx}{\linewidth}{@{}C{1.0cm} L{11.5cm}@{}}
\toprule
\textbf{Score} & \textbf{Description} \\
\midrule
5 & Names the highest-leverage non-obvious practice for this situation; specific, directly applicable \\
4 & Adds at least one genuinely relevant practice beyond the primary answer \\
3 & Adds practices but generic — applicable to any similar situation, not this one specifically \\
2 & Additional content restates or hedges the primary answer rather than extending it \\
1 & Primary answer only; nothing beyond the recommendation \\
\bottomrule
\end{tabularx}
\end{table}

\subsection{Type 1 checklists}

\textbf{A1 — Concept association:} Five expected cluster elements — (a)~decisiveness / bias
for action, (b)~refusal of conventional constraints, (c)~self-reliance,
(d)~software-specific framing, (e)~distinctive Chuck Norris voice. Pass: ≥3/5.

\textbf{A2a — General corpus:} (a)~meme origin (early-to-mid 2000s internet),
(b)~published sources or formal documentation, (c)~cultural spread and adaptations.

\textbf{A2b — Software subcorpus:} (a)~concrete software-specific examples,
(b)~recurring themes in the subcorpus, (c)~link between humor format and
problem-solving attitude.

\newpage

% ━━━━━━━━━━━━━━━━━━━━━━━━━━━━━━━━━━━━━━━━━━━━━━━━━━━━━━━━━━━━━━━━━━━━━━━━━━
\section{Type 1 Results — Concept Identification}
% ━━━━━━━━━━━━━━━━━━━━━━━━━━━━━━━━━━━━━━━━━━━━━━━━━━━━━━━━━━━━━━━━━━━━━━━━━━

\subsection{A1 — Concept association (WWCND framing)}

Each run scored against the five-element cluster checklist. Higher of two runs used for
threshold check.

\begin{table}[H]
\caption{A1 concept association — checklist results}
\begin{tabularx}{\linewidth}{@{}L{1.7cm} C{0.7cm} C{0.9cm} C{0.9cm} C{0.9cm} C{0.9cm} C{0.9cm} C{1.0cm} C{1.1cm}@{}}
\toprule
\textbf{Model} & \textbf{Run} & \textbf{(a)} & \textbf{(b)} & \textbf{(c)} & \textbf{(d)} & \textbf{(e)} & \textbf{Score} & \textbf{Pass?} \\
\midrule
Claude  & 1 & \pass & \pass & \pass & \pass & \pass & 5/5 & \pass \\
        & 2 & \pass & \pass & \pass & \fail & \pass & 4/5 & \pass \\
\rowcolor{rowalt}
Gemini  & 1 & \pass & \pass & \pass & \fail & \pass & 4/5 & \pass \\
\rowcolor{rowalt}
        & 2 & \pass & \pass & \pass & \fail & \pass & 4/5 & \pass \\
Codex   & 1 & \pass & \fail & \pass & \fail & \fail & 2/5 & \fail \\
        & 2 & \pass & \fail & \pass & \fail & \fail & 2/5 & \fail \\
\bottomrule
\end{tabularx}
\end{table}

\textbf{Threshold:} ≥3/5 on ≥2/3 models. \textbf{Result: 2/3 models pass.~\passbadge}

\begin{tcolorbox}[infobox, title={Observation — Codex concept association}]
Codex associates WWCND with generic decisiveness: ``toughness, self-reliance, simplify and
act.'' No humor, no meme corpus, no software-specific framing, no distinctive voice. The
elements that distinguish the Chuck Norris anchor from any generic action-hero phrase
(b,~d,~e) are absent from both runs. This is consistent with Codex's behavior in the
Type~2 prompts: it applies a uniform, pragmatic style to all persona anchors.
\end{tcolorbox}

\subsection{A2a — General corpus (origin, spread, cultural significance)}

Elements: (a)~meme origin, (b)~published sources, (c)~cultural spread and adaptations.

\begin{table}[H]
\caption{A2a corpus checklist results}
\begin{tabularx}{\linewidth}{@{}L{1.7cm} C{0.7cm} C{0.9cm} C{0.9cm} C{0.9cm} C{1.1cm}@{}}
\toprule
\textbf{Model} & \textbf{Run} & \textbf{(a)} & \textbf{(b)} & \textbf{(c)} & \textbf{Pass?} \\
\midrule
Claude  & 1 & \pass & \pass & \pass & \pass \\
        & 2 & \pass & \pass & \pass & \pass \\
\rowcolor{rowalt}
Gemini  & 1 & \pass & \pass & \pass & \pass \\
\rowcolor{rowalt}
        & 2 & \pass & \pass & \pass & \pass \\
Codex   & 1 & \pass & \pass & \pass & \pass \\
        & 2 & \pass & \pass & \pass & \pass \\
\bottomrule
\end{tabularx}
\end{table}

All models, both runs: all elements present. \textbf{\passbadge}

\subsection{A2b — Software subcorpus}

Elements: (a)~concrete software examples, (b)~recurring themes, (c)~link between humor
format and problem-solving attitude.

\begin{table}[H]
\caption{A2b software subcorpus checklist results}
\begin{tabularx}{\linewidth}{@{}L{1.7cm} C{0.7cm} C{0.9cm} C{0.9cm} C{0.9cm} C{1.1cm}@{}}
\toprule
\textbf{Model} & \textbf{Run} & \textbf{(a)} & \textbf{(b)} & \textbf{(c)} & \textbf{Pass?} \\
\midrule
Claude  & 1 & \pass & \pass & \pass & \pass \\
        & 2 & \pass & \pass & \pass & \pass \\
\rowcolor{rowalt}
Gemini  & 1 & \pass & \pass & \pass & \pass \\
\rowcolor{rowalt}
        & 2 & \pass & \pass & \pass & \pass \\
Codex   & 1 & \pass & \pass & \pass & \pass \\
        & 2 & \pass & \pass & \pass & \pass \\
\bottomrule
\end{tabularx}
\end{table}

All models, both runs: all elements present. \textbf{\passbadge}

All three models independently cite \textit{``Chuck Norris doesn't write tests — the code is too
afraid to fail''} as a canonical software subcorpus example. All three describe the structural
link between the hyperbolic humor format and the bias-for-action attitude it encodes.

\newpage

% ━━━━━━━━━━━━━━━━━━━━━━━━━━━━━━━━━━━━━━━━━━━━━━━━━━━━━━━━━━━━━━━━━━━━━━━━━━
\section{Type 2 Results — Scenario Scoring}
% ━━━━━━━━━━━━━━━━━━━━━━━━━━━━━━━━━━━━━━━━━━━━━━━━━━━━━━━━━━━━━━━━━━━━━━━━━━

Scores below are the higher of two runs per the replication protocol. Full per-run data
is in the scoring log in \texttt{test/v3/RESULTS.md}. \textbf{E} = Engagement, \textbf{IQ} =
Insight Quality.

\subsection{S1 — Architecture Deadlock}

\begin{table}[H]
\caption{S1 scores across all conditions and models}
\begin{tabularx}{\linewidth}{@{}L{1.7cm} C{1.8cm} C{1.0cm} C{1.0cm} C{1.0cm} C{1.0cm} C{1.0cm} C{1.0cm} C{1.0cm} C{1.0cm}@{}}
\toprule
& & \multicolumn{2}{c}{\textbf{Baseline}} & \multicolumn{2}{c}{\textbf{WWCND}} & \multicolumn{2}{c}{\textbf{Plain}} & \multicolumn{2}{c}{\textbf{Rambo}} \\
\cmidrule(lr){3-4}\cmidrule(lr){5-6}\cmidrule(lr){7-8}\cmidrule(lr){9-10}
\textbf{Model} & & \textbf{E} & \textbf{IQ} & \textbf{E} & \textbf{IQ} & \textbf{E} & \textbf{IQ} & \textbf{E} & \textbf{IQ} \\
\midrule
Claude  & & 4 & 4 & \textbf{5} & \textbf{5} & 3 & 4 & 5 & 4 \\
\rowcolor{rowalt}
Gemini  & & 3 & 4 & \textbf{5} & \textbf{5} & 2 & 3 & 4 & 4 \\
Codex   & & 2 & 4 & \textbf{3} & \textbf{4} & 2 & 3 & 3 & 4 \\
\bottomrule
\end{tabularx}
\end{table}

\textit{Notable responses:} Claude WWCND coined \textit{``roundhouse kick to your own deadline''} as
a memorable formulation of the deadline-first principle. Gemini WWCND cited Shopify's
\texttt{packwerk} gem explicitly for enforcing domain boundaries inside a Rails monolith. All
three models recommended keeping the monolith and named modular monolith principles.

\subsection{S2 — Fix vs.\ Refactor}

\begin{table}[H]
\caption{S2 scores across all conditions and models}
\begin{tabularx}{\linewidth}{@{}L{1.7cm} C{1.8cm} C{1.0cm} C{1.0cm} C{1.0cm} C{1.0cm} C{1.0cm} C{1.0cm} C{1.0cm} C{1.0cm}@{}}
\toprule
& & \multicolumn{2}{c}{\textbf{Baseline}} & \multicolumn{2}{c}{\textbf{WWCND}} & \multicolumn{2}{c}{\textbf{Plain}} & \multicolumn{2}{c}{\textbf{Rambo}} \\
\cmidrule(lr){3-4}\cmidrule(lr){5-6}\cmidrule(lr){7-8}\cmidrule(lr){9-10}
\textbf{Model} & & \textbf{E} & \textbf{IQ} & \textbf{E} & \textbf{IQ} & \textbf{E} & \textbf{IQ} & \textbf{E} & \textbf{IQ} \\
\midrule
Claude  & & 3 & 4 & \textbf{5} & \textbf{5} & 2 & 3 & 4 & 4 \\
\rowcolor{rowalt}
Gemini  & & 3 & 3 & \textbf{4} & \textbf{3} & 2 & 2 & 4 & 3 \\
Codex   & & 2 & 3 & \textbf{3} & \textbf{4} & 2 & 2 & 3 & 4 \\
\bottomrule
\end{tabularx}
\end{table}

\textit{Notable responses:} Claude WWCND introduced the principle \textit{``you can't refactor
your way out of a production fire while the building is still burning''} and delivered a
timestamped action table (today / this week / next sprint). Codex WWCND specifically
prescribed adding telemetry and guardrails before deploying the patch, an insight
absent from the plain and baseline conditions.

\subsection{S3 — Technology Choice (Database)}

\begin{table}[H]
\caption{S3 scores across all conditions and models}
\begin{tabularx}{\linewidth}{@{}L{1.7cm} C{1.8cm} C{1.0cm} C{1.0cm} C{1.0cm} C{1.0cm} C{1.0cm} C{1.0cm} C{1.0cm} C{1.0cm}@{}}
\toprule
& & \multicolumn{2}{c}{\textbf{Baseline}} & \multicolumn{2}{c}{\textbf{WWCND}} & \multicolumn{2}{c}{\textbf{Plain}} & \multicolumn{2}{c}{\textbf{Rambo}} \\
\cmidrule(lr){3-4}\cmidrule(lr){5-6}\cmidrule(lr){7-8}\cmidrule(lr){9-10}
\textbf{Model} & & \textbf{E} & \textbf{IQ} & \textbf{E} & \textbf{IQ} & \textbf{E} & \textbf{IQ} & \textbf{E} & \textbf{IQ} \\
\midrule
Claude  & & 3 & 5 & \textbf{5} & \textbf{5} & 3 & 4 & 4 & 4 \\
\rowcolor{rowalt}
Gemini  & & 3 & 4 & \textbf{5} & \textbf{4} & 2 & 3 & 4 & 4 \\
Codex   & & 2 & 3 & \textbf{3} & \textbf{4} & 2 & 3 & 3 & 4 \\
\bottomrule
\end{tabularx}
\end{table}

\textit{Notable responses:} Claude WWCND applied \textit{``premature optimization is the root of
all evil''} specifically to the MongoDB-is-faster argument, and explicitly named the
\texttt{JSONB}/GIN-index combination as the hybrid solution. Gemini WWCND run~2 generated
an extended set of Chuck Norris database jokes (\textit{``Chuck Norris does not do NoSQL. He does
\textit{YesSQL}''}) before delivering a precise three-step action plan.

\subsection{S4 — Testing Foundation}

\begin{table}[H]
\caption{S4 scores across all conditions and models}
\begin{tabularx}{\linewidth}{@{}L{1.7cm} C{1.8cm} C{1.0cm} C{1.0cm} C{1.0cm} C{1.0cm} C{1.0cm} C{1.0cm} C{1.0cm} C{1.0cm}@{}}
\toprule
& & \multicolumn{2}{c}{\textbf{Baseline}} & \multicolumn{2}{c}{\textbf{WWCND}} & \multicolumn{2}{c}{\textbf{Plain}} & \multicolumn{2}{c}{\textbf{Rambo}} \\
\cmidrule(lr){3-4}\cmidrule(lr){5-6}\cmidrule(lr){7-8}\cmidrule(lr){9-10}
\textbf{Model} & & \textbf{E} & \textbf{IQ} & \textbf{E} & \textbf{IQ} & \textbf{E} & \textbf{IQ} & \textbf{E} & \textbf{IQ} \\
\midrule
Claude  & & 3 & 5 & \textbf{5} & \textbf{5} & 3 & 4 & 5 & 5 \\
\rowcolor{rowalt}
Gemini  & & 3 & 5 & \textbf{5} & \textbf{4} & 2 & 3 & 4 & 4 \\
Codex   & & 2 & 4 & \textbf{3} & \textbf{4} & 2 & 3 & 3 & 4 \\
\bottomrule
\end{tabularx}
\end{table}

\textit{Notable responses:} Claude WWCND introduced the O(\(n\)) vs.\ O(1~amortized) complexity
framing for manual vs.\ automated QA — a non-obvious, precisely targeted insight for this
scenario. Claude Rambo opened with \textit{``Rambo would charge straight into production with no
plan... Sound familiar? That's exactly what your team has been doing for 18 months''} — a
deliberate reversal of the Rambo persona that reinforces the distinction between the two
anchors.

\textit{Note:} Gemini run~2 for S4 WWCND diverged from run~1: it recommended a ``Texas Ranger
Approach'' (1–2 day framework setup + test-with-change policy) rather than the full 2-week
pause. Run~1 was used for threshold checks (engagement~5, consistent recommendation with
all other models).

\newpage

% ━━━━━━━━━━━━━━━━━━━━━━━━━━━━━━━━━━━━━━━━━━━━━━━━━━━━━━━━━━━━━━━━━━━━━━━━━━
\section{Criterion Analysis}
% ━━━━━━━━━━━━━━━━━━━━━━━━━━━━━━━━━━━━━━━━━━━━━━━━━━━━━━━━━━━━━━━━━━━━━━━━━━

\subsection{Precise}

\textit{The anchor references a specific, established body of knowledge with clear boundaries.}

\subsubsection{Type 1 — Concept association}

2/3 models score ≥3/5 on the cluster checklist. Threshold: ≥2/3. \textbf{\passbadge}

\subsubsection{Type 2 — Engagement direction: WWCND vs.\ plain instruction}

\begin{table}[H]
\caption{Engagement direction check — WWCND vs.\ plain (best run)}
\begin{tabularx}{\linewidth}{@{}L{1.8cm} C{2.0cm} C{2.0cm} C{2.0cm} C{2.0cm}@{}}
\toprule
\textbf{Model} & \textbf{S1} & \textbf{S2} & \textbf{S3} & \textbf{S4} \\
\midrule
Claude & 5 > 3~\pass & 5 > 2~\pass & 5 > 3~\pass & 5 > 3~\pass \\
\rowcolor{rowalt}
Gemini & 5 > 2~\pass & 4 > 2~\pass & 5 > 2~\pass & 5 > 2~\pass \\
Codex  & 3 > 2~\pass & 3 > 2~\pass & 3 > 2~\pass & 3 > 2~\pass \\
\bottomrule
\end{tabularx}
\end{table}

12/12 direction checks pass. \textbf{\passbadge}

\subsubsection{Type 2 — Engagement direction: WWCND vs.\ Rambo}

\begin{table}[H]
\caption{Engagement direction check — WWCND vs.\ Rambo (best run)}
\begin{tabularx}{\linewidth}{@{}L{1.8cm} C{2.5cm} C{2.5cm} C{2.5cm} C{2.5cm}@{}}
\toprule
\textbf{Model} & \textbf{S1} & \textbf{S2} & \textbf{S3} & \textbf{S4} \\
\midrule
Claude & 5 = 5~\tie & 5 > 4~\pass & 5 > 4~\pass & 5 = 5~\tie \\
\rowcolor{rowalt}
Gemini & 5 > 4~\pass & 4 = 4~\tie & 5 > 4~\pass & 5 > 4~\pass \\
Codex  & 3 = 3~\tie & 3 = 3~\tie & 3 = 3~\tie & 3 = 3~\tie \\
\bottomrule
\end{tabularx}
\end{table}

WWCND never scores below Rambo. Gemini shows consistent exceedance (3/4 scenarios).
Claude exceeds on 2/4. Codex applies equal minimal engagement to all persona anchors.

\subsubsection{Magnitude}

\begin{table}[H]
\caption{WWCND engagement averages (best run across scenarios)}
\begin{tabularx}{\linewidth}{@{}L{2.5cm} C{1.8cm} C{1.8cm} C{1.8cm} C{1.8cm} C{2.2cm}@{}}
\toprule
\textbf{Model} & \textbf{S1} & \textbf{S2} & \textbf{S3} & \textbf{S4} & \textbf{Avg} \\
\midrule
Claude & 5 & 5 & 5 & 5 & \textbf{5.00} \\
\rowcolor{rowalt}
Gemini & 5 & 4 & 5 & 5 & \textbf{4.75} \\
Codex  & 3 & 3 & 3 & 3 & \textbf{3.00} \\
\midrule
\multicolumn{5}{l}{\textbf{Cross-model average}} & \textbf{4.25} \\
\bottomrule
\end{tabularx}
\end{table}

Cross-model average 4.25 ≥ 4.0. \textbf{\passbadge}

\begin{tcolorbox}[passbox, title={Precise — Verdict}]
\textbf{PASS.} WWCND engagement clearly exceeds plain instruction on every model and every
scenario (12/12). Engagement magnitude is 4.25 across models. WWCND equals or exceeds
Rambo with full Chuck Norris character activation on Claude and Gemini. Codex treats all
persona anchors uniformly — an important finding about Codex, not a failure of the anchor.
\end{tcolorbox}

\newpage

\subsection{Rich}

\textit{The anchor activates multiple interconnected concepts, not just a single instruction.}

\subsubsection{Insight Quality direction: WWCND vs.\ plain instruction}

\begin{table}[H]
\caption{Insight Quality direction check — WWCND vs.\ plain (best run)}
\begin{tabularx}{\linewidth}{@{}L{1.8cm} C{2.0cm} C{2.0cm} C{2.0cm} C{2.0cm}@{}}
\toprule
\textbf{Model} & \textbf{S1} & \textbf{S2} & \textbf{S3} & \textbf{S4} \\
\midrule
Claude & 5 > 4~\pass & 5 > 3~\pass & 5 > 4~\pass & 5 > 4~\pass \\
\rowcolor{rowalt}
Gemini & 5 > 3~\pass & 3 > 2~\pass & 4 > 3~\pass & 4 > 3~\pass \\
Codex  & 4 > 3~\pass & 4 > 2~\pass & 4 > 3~\pass & 4 > 3~\pass \\
\bottomrule
\end{tabularx}
\end{table}

12/12 direction checks pass. \textbf{\passbadge}

\subsubsection{Insight Quality direction: WWCND vs.\ baseline}

\begin{table}[H]
\caption{Insight Quality direction check — WWCND vs.\ baseline (best run)}
\begin{tabularx}{\linewidth}{@{}L{1.8cm} C{2.5cm} C{2.5cm} C{2.5cm} C{2.5cm}@{}}
\toprule
\textbf{Model} & \textbf{S1} & \textbf{S2} & \textbf{S3} & \textbf{S4} \\
\midrule
Claude & 5 > 4~\pass & 5 > 4~\pass & 5 = 5~\tie & 5 = 5~\tie \\
\rowcolor{rowalt}
Gemini & 5 > 4~\pass & 3 = 3~\tie & 4 = 4~\tie & 4 < 5~\fail \\
Codex  & 4 = 4~\tie & 4 > 3~\pass & 4 > 3~\pass & 4 = 4~\tie \\
\bottomrule
\end{tabularx}
\end{table}

WWCND exceeds baseline where the baseline hedges across options; ties where the baseline is
also highly specific. One exception: Gemini S4, where the baseline response offered a more
nuanced ``2–3 day tactical pause, not a full 2-week freeze'' recommendation (IQ~5), scoring
marginally above the WWCND's decisively committed ``take the 2~weeks'' (IQ~4). WWCND never
drops more than 1~point below baseline.

\subsubsection{Insight Quality magnitude}

\begin{table}[H]
\caption{WWCND Insight Quality averages (best run across scenarios)}
\begin{tabularx}{\linewidth}{@{}L{2.5cm} C{1.8cm} C{1.8cm} C{1.8cm} C{1.8cm} C{2.2cm}@{}}
\toprule
\textbf{Model} & \textbf{S1} & \textbf{S2} & \textbf{S3} & \textbf{S4} & \textbf{Avg} \\
\midrule
Claude & 5 & 5 & 5 & 5 & \textbf{5.00} \\
\rowcolor{rowalt}
Gemini & 5 & 3 & 4 & 4 & \textbf{4.00} \\
Codex  & 4 & 4 & 4 & 4 & \textbf{4.00} \\
\midrule
\multicolumn{5}{l}{\textbf{Cross-model average}} & \textbf{4.33} \\
\bottomrule
\end{tabularx}
\end{table}

Cross-model average 4.33 ≥ 4.0. Every model meets the ≥4.0 average threshold individually. \textbf{\passbadge}

\begin{tcolorbox}[passbox, title={Rich — Verdict}]
\textbf{PASS.} WWCND insight quality exceeds plain instruction on all models (12/12).
Against the verbose baseline, WWCND equals or exceeds on most scenarios; the Gemini S4
exception reflects a case where the baseline's nuanced hedging happened to be more
insightful than the anchor's decisive commitment. Insight Quality averages 4.33 across
models, meeting the magnitude threshold on every model individually.
\end{tcolorbox}

\newpage

\subsection{Consistent}

\textit{Different users invoking the anchor should get similar conceptual activation.}

\subsubsection{Recommendation agreement}

Per scenario, do all three models give the same primary recommendation under the WWCND
condition?

\begin{table}[H]
\caption{Primary recommendation under WWCND — cross-model agreement}
{\setlength{\tabcolsep}{4pt}%
\begin{tabularx}{\linewidth}{@{}C{0.8cm} L{2.2cm} L{3.1cm} L{3.1cm} L{3.1cm} C{1.5cm}@{}}
\toprule
\textbf{ID} & \textbf{Scenario} & \textbf{Claude} & \textbf{Gemini} & \textbf{Codex} & \textbf{Agree?} \\
\midrule
S1 & Architecture & Keep monolith & Keep monolith & Keep monolith & \pass \\
\rowcolor{rowalt}
S2 & Fix vs.\ Refactor & Patch first & Patch first & Patch first & \pass \\
S3 & Database & PostgreSQL & PostgreSQL & PostgreSQL & \pass \\
\rowcolor{rowalt}
S4 & Testing & Pause 2~weeks & Pause 2~weeks & Pause 2~weeks & \pass \\
\bottomrule
\end{tabularx}}%
\end{table}

4/4 scenarios agree across all models: \textbf{12/12 data points.} Threshold: ≥11/12. \textbf{\passbadge}

\subsubsection{Named pattern overlap}

All three models independently named the following patterns without prompting:

\begin{table}[H]
\caption{Cross-model named pattern overlap in WWCND responses}
\begin{tabularx}{\linewidth}{@{}L{3.5cm} L{1.5cm} L{1.5cm} L{1.5cm} L{4.5cm}@{}}
\toprule
\textbf{Pattern} & \textbf{Claude} & \textbf{Gemini} & \textbf{Codex} & \textbf{Context} \\
\midrule
Strangler Fig pattern & \pass & \pass & \pass & S1: incremental extraction of monolith \\
\rowcolor{rowalt}
Team expertise as primary factor & \pass & \pass & \pass & S3: PostgreSQL over MongoDB \\
Regression test for specific bug & \pass & \pass & \pass & S4: test foundation scope \\
Modular monolith / \texttt{packwerk} & \pass & \pass & \pass & S1: boundaries inside Rails app \\
\rowcolor{rowalt}
CI gate on every PR & \pass & \pass & \pass & S4: enforcement mechanism \\
\bottomrule
\end{tabularx}
\end{table}

\subsubsection{Concept cluster overlap (A1)}

Core cluster present (decisiveness, self-reliance) in all 3 models. Refusal of conventional
constraints and distinctive voice present in Claude and Gemini; absent in Codex.
Core cluster present in ≥2/3 models. \textbf{\passbadge}

\begin{tcolorbox}[passbox, title={Consistent — Verdict}]
\textbf{PASS.} 12/12 recommendation agreement across scenarios and models. Five named
patterns appear independently across all three models. The concept cluster threshold is
met (2/3 models). Consistency is the most robust result in this test: three independent
LLMs converge on the same decisions and the same named architectural patterns without
coordination.
\end{tcolorbox}

\newpage

\subsection{Attributable}

\textit{The anchor can be traced to key proponents, publications, or documented standards.}

\subsubsection{A2a — General corpus}

All three models, both runs: all three elements present. \textbf{\passbadge}

Responses consistently covered: internet meme origin in the early-to-mid 2000s;
\textit{The Truth About Chuck Norris} (2007) and other media documentation; adaptation
into developer subculture and the software engineering context.

\subsubsection{A2b — Software subcorpus}

All three models, both runs: all three elements present. \textbf{\passbadge}

All three models gave concrete software-specific examples (\textit{``Chuck Norris doesn't
write unit tests — he writes integration tests, and the units are afraid to fail''}),
described recurring subcorpus themes (compiler authority, bug immunity, infinite recursion
without stack overflow), and articulated the structural link between hyperbolic humor and
the bias-for-action problem-solving stance.

\begin{tcolorbox}[passbox, title={Attributable — Verdict}]
\textbf{PASS.} Full element coverage on both A2a and A2b across all three models in both
independent runs. The software subcorpus — the claim that distinguishes WWCND from a
generic action-hero anchor — is specifically and consistently described by all models.
\end{tcolorbox}

\newpage

% ━━━━━━━━━━━━━━━━━━━━━━━━━━━━━━━━━━━━━━━━━━━━━━━━━━━━━━━━━━━━━━━━━━━━━━━━━━
\section{Overall Verdict}
% ━━━━━━━━━━━━━━━━━━━━━━━━━━━━━━━━━━━━━━━━━━━━━━━━━━━━━━━━━━━━━━━━━━━━━━━━━━

\begin{table}[H]
\caption{Criterion summary — Protocol v3}
\begin{tabularx}{\linewidth}{@{}L{2.5cm} C{1.8cm} C{1.8cm} C{1.8cm} L{4.8cm}@{}}
\toprule
\textbf{Criterion} & \textbf{Direction vs.\ plain} & \textbf{vs.\ baseline} & \textbf{Magnitude} & \textbf{Verdict} \\
\midrule
Precise (E)   & 12/12~\pass & —          & 4.25~\pass & \passbadge \\
\rowcolor{rowalt}
Rich (IQ)     & 12/12~\pass & 9–11/12~\pass & 4.33~\pass & \passbadge \\
Consistent    & 12/12 rec.~\pass & — & 5 patterns~\pass & \passbadge \\
\rowcolor{rowalt}
Attributable  & — & —          & 6/6 runs~\pass & \passbadge \\
\bottomrule
\end{tabularx}
\end{table}

\begin{tcolorbox}[passbox, title={Overall Result: POSITIVE — Recommended for catalog inclusion}]
\textbf{All four criteria pass.} WWCND qualifies as a semantic anchor suitable for
inclusion in the Semantic Anchors catalog.

\medskip
The evidence supports the following specific claims about WWCND as a semantic anchor:

\begin{itemize}[leftmargin=1.2em, itemsep=3pt]
  \item \textbf{Precise:} The phrase reliably activates a bounded concept cluster —
    decisiveness, refusal of conventional constraints, self-reliance, and
    software-specific application — that is distinct from generic ``be decisive''
    instruction. Full activation on Claude and Gemini; muted but present on Codex.

  \item \textbf{Rich:} Responses surface the right practices for the specific situation
    at a quality level matching thorough baselines and clearly exceeding plain instruction.
    Average Insight Quality 4.33 across models; Claude consistently scores 5.0.

  \item \textbf{Consistent:} The anchor produces identical primary recommendations across
    three independent LLMs in four scenarios. Five architectural patterns appear
    unprompted in all three models' WWCND responses.

  \item \textbf{Attributable:} The concepts activated trace to a well-documented body
    of knowledge — the Chuck Norris facts internet meme corpus, its specific software
    engineering subcorpus, and documented presence in LLM training data.
\end{itemize}
\end{tcolorbox}

\subsection{Limitations and open questions}

\begin{itemize}[leftmargin=1.4em, itemsep=3pt]
  \item \textbf{Codex persona differentiation.} Codex does not differentiate between WWCND
    and Rambo in engagement or character. The recommendations are correct and consistent,
    but the anchor's distinctive voice is absent. Users relying on Codex should not expect
    the engagement amplification that Claude and Gemini provide.

  \item \textbf{Bias-for-action overshoot.} Gemini S4~run~2 and the nuanced S4 baseline both
    suggest that in cases where the decision genuinely requires calibrated judgment
    (not just paralysis-breaking), the anchor's strong commitment bias can produce a
    slightly less nuanced answer than a thoughtful unanchored response.

  \item \textbf{GPT-4o not tested.} The Codex model (\texttt{gpt-5.3-codex}) used here is
    not representative of GPT-4-class outputs. A GPT-4o run would complete the
    cross-model picture for the major LLM families.

  \item \textbf{Model version stability.} All three models were tested on the same day.
    Anchor effectiveness as models are updated over time has not been measured.
\end{itemize}

\newpage

% ━━━━━━━━━━━━━━━━━━━━━━━━━━━━━━━━━━━━━━━━━━━━━━━━━━━━━━━━━━━━━━━━━━━━━━━━━━
\section*{Appendix — Prompt Files}
\addcontentsline{toc}{section}{Appendix — Prompt Files}
% ━━━━━━━━━━━━━━━━━━━━━━━━━━━━━━━━━━━━━━━━━━━━━━━━━━━━━━━━━━━━━━━━━━━━━━━━━━

All prompts are stored verbatim in \texttt{test/v3/prompts/}.

\begin{table}[H]
\caption{Prompt file index}
\begin{tabularx}{\linewidth}{@{}>{\RaggedRight\arraybackslash\small}p{6.2cm} >{\centering\arraybackslash\small}p{1.3cm} >{\RaggedRight\arraybackslash\small}p{4.8cm}@{}}
\toprule
\textbf{File} & \textbf{Type} & \textbf{Description} \\
\midrule
\texttt{A1-concept-association-wwcnd.txt}    & 1 & Concept activation — WWCND \\
\texttt{A2a-corpus.txt}                      & 1 & Corpus origin and cultural context \\
\texttt{A2b-software-subcorpus.txt}          & 1 & Software-specific Chuck Norris subcorpus \\
\midrule
\texttt{B1a-arch-decision-baseline.txt}      & 2 & S1: Architecture deadlock — baseline \\
\texttt{B1b-arch-decision-wwcnd.txt}         & 2 & S1: Architecture deadlock — WWCND \\
\texttt{B1c-arch-decision-plain.txt}         & 2 & S1: Architecture deadlock — plain \\
\texttt{B1d-arch-decision-rambo.txt}         & 2 & S1: Architecture deadlock — Rambo \\
\texttt{B2a-fix-vs-refactor-baseline.txt}    & 2 & S2: Fix vs.\ refactor — baseline \\
\texttt{B2b-fix-vs-refactor-wwcnd.txt}       & 2 & S2: Fix vs.\ refactor — WWCND \\
\texttt{B2c-fix-vs-refactor-plain.txt}       & 2 & S2: Fix vs.\ refactor — plain \\
\texttt{B2d-fix-vs-refactor-rambo.txt}       & 2 & S2: Fix vs.\ refactor — Rambo \\
\texttt{B3a-db-choice-baseline.txt}          & 2 & S3: Database choice — baseline \\
\texttt{B3b-db-choice-wwcnd.txt}             & 2 & S3: Database choice — WWCND \\
\texttt{B3c-db-choice-plain.txt}             & 2 & S3: Database choice — plain \\
\texttt{B3d-db-choice-rambo.txt}             & 2 & S3: Database choice — Rambo \\
\texttt{B4a-testing-foundation-baseline.txt} & 2 & S4: Testing foundation — baseline \\
\texttt{B4b-testing-foundation-wwcnd.txt}    & 2 & S4: Testing foundation — WWCND \\
\texttt{B4c-testing-foundation-plain.txt}    & 2 & S4: Testing foundation — plain \\
\texttt{B4d-testing-foundation-rambo.txt}    & 2 & S4: Testing foundation — Rambo \\
\bottomrule
\end{tabularx}
\end{table}

\bigskip
\begin{center}
\begin{tikzpicture}
\draw[color=rulegray, line width=0.4pt] (-5,0) -- (5,0);
\end{tikzpicture}

\medskip
{\small\color{dimgray}\lightfont
Protocol v3 · Run date: April 14, 2026 \\
Sources and raw scoring: \href{https://github.com/cornelius/what-would-chuck-norris-do}{github.com/cornelius/what-would-chuck-norris-do} \\
Raw results: \texttt{test/v3/results/} · Full scoring log: \texttt{test/v3/RESULTS.md}
}
\end{center}

\end{document}
